\documentclass[12pt]{article}
\usepackage[ukrainian,shorthands=off]{babel}   % shorthands=off: символ " не активний (потрібно для міток "..." у TikZ всередині \zadtask)
\usepackage{fontspec}
% --- НАЛАШТУВАННЯ ШРИФТІВ ---
% Шрифти Schoolbook-*.otf лежать у корені проєкту. Шлях підбирається автоматично:
% ../ (компіляція з папки теми), ./ (корінь проєкту), ../../ (підпапки теми 28); інакше -- TeX Gyre Schola.
\newcommand{\nmtSchoolbook}[1]{\setmainfont{Schoolbook-Regular.otf}[
    Path = #1,
    BoldFont = Schoolbook-Bold.otf,
    ItalicFont = Schoolbook-Italic.otf,
    BoldItalicFont = Schoolbook-BoldItalic.otf
]}
\IfFileExists{../Schoolbook-Regular.otf}{\nmtSchoolbook{../}}{%
 \IfFileExists{./Schoolbook-Regular.otf}{\nmtSchoolbook{./}}{%
  \IfFileExists{../../Schoolbook-Regular.otf}{\nmtSchoolbook{../../}}{%
   \setmainfont{texgyreschola}[Extension=.otf, UprightFont=*-regular, BoldFont=*-bold, ItalicFont=*-italic, BoldItalicFont=*-bolditalic]}}}
\usepackage[italic]{mathastext}
\usepackage[a4paper,margin=1.8cm,bottom=2cm,top=2cm]{geometry}
\usepackage{amsmath,amssymb}
\usepackage{tikz}
\usepackage{pgfplots}
\pgfplotsset{compat=1.18}
\usepgfplotslibrary{fillbetween}
\usetikzlibrary{calc,patterns,angles,quotes,intersections,babel,3d,shapes.symbols,shapes.geometric,decorations.pathreplacing,shadings,arrows.meta,hobby}
\usepackage{xcolor,array,fancyhdr,enumitem,multicol}
\usepackage{colortbl,multirow,diagbox,needspace}
\usepackage{tcolorbox}
\tcbuselibrary{skins,breakable}
\usepackage[hidelinks]{hyperref}
% водяний знак; щоб зібрати без нього (версія для вчителів), перед \documentclass пишуть \def\nmtnowatermark{}
\ifdefined\nmtnowatermark\else
\usepackage[firstpage=false, text=@pvtr2525, scale=0.6, color=gray!12]{draftwatermark}
\fi

% --- АВТОСТИСКАННЯ: рисунок/сітка, ширша за колонку, масштабується до ширини колонки ---
\newsavebox{\nmtfitbox}
\newenvironment{nmtfit}{\begin{lrbox}{\nmtfitbox}}{\end{lrbox}%
  \ifdim\wd\nmtfitbox>\linewidth \resizebox{\linewidth}{!}{\usebox{\nmtfitbox}}\else\usebox{\nmtfitbox}\fi}
\let\nmtIncludegraphicsOrig\includegraphics
\renewcommand{\includegraphics}[2][]{\begin{nmtfit}\nmtIncludegraphicsOrig[#1]{#2}\end{nmtfit}}


\definecolor{mainGreen}{RGB}{34, 120, 64}
\definecolor{yearOrange}{RGB}{220, 100, 30}
\definecolor{yearcolor}{RGB}{220, 100, 30}
\definecolor{titleBg}{RGB}{225, 240, 220}
\definecolor{instrBg}{RGB}{255, 240, 220}
\definecolor{instrBorder}{RGB}{220, 150, 80}
\definecolor{headerblue}{RGB}{0, 102, 204}   % використовується в рисунках старої бази

\setlength{\parindent}{0pt}
\emergencystretch=1.5em

\pagestyle{fancy}
\fancyhf{}
\renewcommand{\headrulewidth}{0.4pt}
\renewcommand{\footrulewidth}{0.4pt}
\fancyhead[C]{\small\color{gray!80} НМТ 2026, 8 червня}
\fancyhead[R]{\small\color{mainGreen}\textbf{\href{https://t.me/pvtr2525}{@pvtr2525}}}
\fancyfoot[L]{\small\color{gray!80}\href{https://t.me/pvtr2525}{ТГ: @pvtr2525}}
\fancyfoot[C]{\small\color{gray!80}\thepage}
\fancyfoot[R]{\small\color{gray!80}НМТ з математики}

% --- ТАБЛИЦІ ВІДПОВІДЕЙ ---
% ширина клітинки = (ширина рядка - відступи) / 5: на всю сторінку це ~3 см, у minipage поруч із рисунком таблиця стискається;
% п'ять колонок |m{w}| мають 10 відступів \tabcolsep і 6 вертикальних ліній -- тоді таблиця рівно на \linewidth
\newlength{\nmtcellw}
\newcommand{\nmtSetCell}{\setlength{\nmtcellw}{\dimexpr(\linewidth-10\tabcolsep-6\arrayrulewidth)/5\relax}}
\newcommand{\answerTable}[5]{%
\par\nopagebreak\vspace{0.15cm}
\noindent\nmtSetCell
\begin{tabular}{|*{5}{>{\centering\arraybackslash}m{\nmtcellw}|}}
\hline
\rule[-0.15cm]{0pt}{0.6cm}\textbf{А} & \rule[-0.15cm]{0pt}{0.6cm}\textbf{Б} & \rule[-0.15cm]{0pt}{0.6cm}\textbf{В} & \rule[-0.15cm]{0pt}{0.6cm}\textbf{Г} & \rule[-0.15cm]{0pt}{0.6cm}\textbf{Д} \\
\hline
\rule[-0.2cm]{0pt}{0.75cm}#1 & \rule[-0.2cm]{0pt}{0.75cm}#2 & \rule[-0.2cm]{0pt}{0.75cm}#3 & \rule[-0.2cm]{0pt}{0.75cm}#4 & \rule[-0.2cm]{0pt}{0.75cm}#5 \\
\hline
\end{tabular}
\par\vspace{0.25cm}
}
\newcommand{\answerTableTall}[5]{%
\par\nopagebreak\vspace{0.15cm}
\noindent\nmtSetCell
\begin{tabular}{|*{5}{>{\centering\arraybackslash}m{\nmtcellw}|}}
\hline
\rule[-0.2cm]{0pt}{0.7cm}\textbf{А} & \rule[-0.2cm]{0pt}{0.7cm}\textbf{Б} & \rule[-0.2cm]{0pt}{0.7cm}\textbf{В} & \rule[-0.2cm]{0pt}{0.7cm}\textbf{Г} & \rule[-0.2cm]{0pt}{0.7cm}\textbf{Д} \\
\hline
\rule[-0.45cm]{0pt}{1.2cm}#1 & \rule[-0.45cm]{0pt}{1.2cm}#2 & \rule[-0.45cm]{0pt}{1.2cm}#3 & \rule[-0.45cm]{0pt}{1.2cm}#4 & \rule[-0.45cm]{0pt}{1.2cm}#5 \\
\hline
\end{tabular}
\par\vspace{0.25cm}
}
\newcommand{\instructionBox}[1]{%
\par\vspace{0.3cm}
\begin{tcolorbox}[colback=instrBg, colframe=instrBorder, boxrule=0.6pt, arc=2pt,
    left=8pt, right=8pt, top=4pt, bottom=4pt]
\centering\small\bfseries #1
\end{tcolorbox}
\par\vspace{0.15cm}
}
\newcommand{\sectionTitle}[1]{%
\par\vspace{0.4cm}
\begin{tcolorbox}[colback=titleBg, colframe=titleBg, boxrule=0pt, arc=8pt,
    halign=center, fontupper=\Large\bfseries\color{mainGreen}]
\hyphenpenalty=10000 \exhyphenpenalty=10000 #1
\end{tcolorbox}
\par\vspace{0.3cm}
}
\newcommand{\task}[2]{\noindent\textbf{#1.}\ \ #2\par\vspace{0.2cm}}
% --- НМТ-СТИЛЬ ПОЛЕ ВІДПОВІДІ ---
\newcommand{\nmtAnswerBox}{%
\par\nopagebreak\vspace{0.2cm}\noindent
Відповідь:\ \
\begin{tabular}{@{}*{4}{|>{\centering\arraybackslash}p{0.8cm}}|@{\hspace{0.1cm}\textbf{,}\hspace{0.1cm}}*{3}{|>{\centering\arraybackslash}p{0.8cm}}|@{}}
\hline
\rule[-0.2cm]{0pt}{0.7cm} & & & & & & \\
\hline
\end{tabular}
\par\vspace{0.3cm}
}
% --- СІТКА ДЛЯ МАТЧИНГУ ---
\newsavebox{\nmtgridbox}
\newcommand{\matchingGrid}{%
    \sbox{\nmtgridbox}{\matchingGridRaw}%
    \ifdim\wd\nmtgridbox>\linewidth \resizebox{\linewidth}{!}{\usebox{\nmtgridbox}}\else\usebox{\nmtgridbox}\fi}
\newcommand{\matchingGridRaw}{%
    \begingroup
    \renewcommand{\arraystretch}{1.15}
    \setlength{\tabcolsep}{4pt}
    \footnotesize
    \begin{tabular}{r|c|c|c|c|c|}
         \multicolumn{1}{c}{} & \multicolumn{1}{c}{\textbf{А}} & \multicolumn{1}{c}{\textbf{Б}} & \multicolumn{1}{c}{\textbf{В}} & \multicolumn{1}{c}{\textbf{Г}} & \multicolumn{1}{c}{\textbf{Д}} \\ \cline{2-6}
         \textbf{1} & \phantom{X} & \phantom{X} & \phantom{X} & \phantom{X} & \phantom{X} \\ \cline{2-6}
         \textbf{2} & & & & & \\ \cline{2-6}
         \textbf{3} & & & & & \\ \cline{2-6}
    \end{tabular}
    \endgroup
}
% --- ДОДАТКОВІ МАКРОСИ ЗБІРНИКА ---
\newcommand{\taskBlock}[1]{%
\par\vspace{0.45cm}
\noindent{\color{mainGreen}\rule{\linewidth}{0.8pt}}\par\vspace{0.12cm}
\noindent{\small\bfseries\color{mainGreen}#1}\par\vspace{0.25cm}
}
\newcounter{zad}
\newcommand{\zadtask}[1]{\stepcounter{zad}\noindent\textbf{\thezad.}\ \ #1\par\nopagebreak\vspace{0.2cm}}
\newcommand{\zadnum}{\stepcounter{zad}\textbf{\thezad.}\ \ }
\newcounter{ans}
\newcommand{\ansitem}[1]{\stepcounter{ans}\noindent\textbf{\theans.}\ #1\par\vspace{0.07cm}}
% мітка року: нерозривна, притиснута праворуч; якщо не вміщається -- праворуч на новому рядку
\newcommand{\nmtyear}[1]{\unskip\nobreak\finalhyphendemerits=0 \hfill\penalty100\hskip1em\hbox{}\nobreak\hfill{\small\color{yearOrange}\mbox{\rule{0pt}{2.3ex}(НМТ~#1)}}}
% --- МАКРОС КОРОТКОГО РОЗВ'ЯЗКУ ---
\newcommand{\solution}[1]{%
\par\vspace{0.12cm}
\begin{tcolorbox}[colback=titleBg!45, colframe=mainGreen!55, boxrule=0.4pt, arc=2pt,
    left=6pt, right=6pt, top=3pt, bottom=3pt, breakable]
\footnotesize\textbf{Короткий розв'язок.}\ #1
\end{tcolorbox}
\par\vspace{0.12cm}
}
% Розв'язки й відповіді (є до завдань НМТ-2026): \showsolutionstrue -- показати, \showsolutionsfalse -- сховати
\newif\ifshowsolutions
\showsolutionsfalse
% --- ЗАГОЛОВКИ З ПУНКТАМИ КЛІКАБЕЛЬНОГО ЗМІСТУ ---
\setcounter{tocdepth}{2}
\newcommand{\chapterTitle}[1]{%
\clearpage\phantomsection\addcontentsline{toc}{section}{#1}%
\sectionTitle{#1}}
\newcommand{\typeTitle}[1]{%
\needspace{6\baselineskip}%
\phantomsection\addcontentsline{toc}{subsection}{\quad #1}%
\par\vspace{0.3cm}
\noindent{\large\bfseries\color{yearOrange}#1}\par\vspace{0.08cm}
\noindent{\color{yearOrange}\rule{\linewidth}{0.4pt}}\par\nopagebreak\vspace{0.2cm}}
\raggedbottom
% усередині одного завдання сторінка не розривається: нерозривний вертикальний проміжок
% і нерозривна межа абзаців (умова ніколи не відділяється від варіантів, рисунка чи сітки)
\newcommand{\nbvspace}[1]{\nopagebreak\vspace{#1}}
\newcommand{\nmtnobreak}{\ifvmode\penalty10000 \fi}
% --- ЄДИНИЙ МАКЕТ ЗАВДАНЬ НА ВІДПОВІДНІСТЬ ---
% мітка в колонці сталої ширини, текст -- у parbox: продовження довгого рядка
% вирівнюється під текстом, а не під міткою; відступ між пунктами однаковий скрізь
\newlength{\nmtMatchLab}\setlength{\nmtMatchLab}{1.6em}
\newlength{\nmtMatchSep}\setlength{\nmtMatchSep}{0.28cm}
\newcommand{\matchHead}[1]{\par\noindent\textit{#1}\par\nopagebreak\vspace{0.2cm}}
% шаблонний заголовок стовпця зі старого макета -- лише коли завдання не має власного \matchHead
\makeatletter\newcommand{\nmtHeadUnless}[2]{\in@{\matchHead}{#1}\ifin@\else#2\fi}\makeatother
\newcommand{\matchItem}[2]{\par\noindent\makebox[\nmtMatchLab][l]{\textbf{#1}}%
\parbox[t]{\dimexpr\linewidth-\nmtMatchLab\relax}{\raggedright #2}\par\nopagebreak\vspace{\nmtMatchSep}}
\newcommand{\ansTheme}[1]{\par\vspace{0.25cm}\noindent{\bfseries\color{mainGreen}#1}\par\vspace{0.12cm}}
\newcommand{\ansType}[1]{\par\vspace{0.1cm}\noindent{\itshape\color{yearOrange}#1}\par\vspace{0.08cm}}

\graphicspath{{../1. Числа. Звичайні та десяткові дроби. Модуль числа/}{../10. Основи планіметрії/}{../11. Трикутники/}{../12. Паралелограм/}{../13. Прямокутник/}{../14. Квадрат/}{../15. Ромб/}{../16. Трапеція/}{../17. Коло, круг і їх елементи/}{../18. Координати і вектори на площині і в просторі. Рівняння кола/}{../19. Лінійні та квадратні нерівності. Системи лінійних нерівностей. Метод інтервалів/}{../2.відсотки, арифметичні задачі, подільність/}{../20. Функція та її властивості/}{../21.  Графіки елементарних функцій, їх побудова графіків функцій та їх перетворення/}{../22. Модульні  рівняння та нерівності/}{../23. Ірраціональні рівняння та нерівності/}{../24. Арифметична прогресія/}{../25. Геометрична прогресія/}{../26. Тригонометричні вирази/}{../27. Тригонометричні рівняння/}{../28. Показникова функція. Показникові рівняння/Показникові рівняння/}{../28. Показникова функція. Показникові рівняння/Показникові функція і вирази/}{../29. Показникові нерівності/}{../3. Степені. Одночлени. стандартний вигляд числа/}{../30. Логарифм. Логарифмічна функція/}{../31. Логарифмічні рівняння/}{../32. Логарифмічні нерівності/}{../33. Похідна/}{../34. Первісна/}{../35. Комбінаторика/}{../36. Теорія ймовірності/}{../37. Статистика/}{../38. Призма. Паралелепіпед. Куб/}{../39. Піраміда/}{../4.Многочлени. ФСМ/}{../40. Циліндр/}{../41. Конус/}{../42. Куля. Сфера/}{../43. Параметри/}{../5. дробові раціональні вирази і перетворення/}{../6. Лінійні рівняння/}{../7. Квадратні корені. Корені вищих степенів. Раціональний степінь/}{../8. квадратні рівняння, і рівняння, що до них зводяться/}{../9. Системи рівнянь/}}
\renewcommand{\instructionBox}[1]{%
\par\vspace{0.3cm}\noindent
\begin{tcolorbox}[nobeforeafter, width=\linewidth, colback=instrBg, colframe=instrBorder, boxrule=0.6pt, arc=2pt,
    left=8pt, right=8pt, top=4pt, bottom=4pt]
\centering\small\bfseries #1
\end{tcolorbox}
\par\nopagebreak\vspace{0.15cm}\nopagebreak
}
\renewcommand{\nmtAnswerBox}{%
\par\nopagebreak\vspace{0.2cm}\noindent
Відповідь:\ \
\begin{tabular}{@{}*{5}{|>{\centering\arraybackslash}p{0.8cm}}|@{\hspace{0.1cm}\textbf{,}\hspace{0.1cm}}*{3}{|>{\centering\arraybackslash}p{0.8cm}}|@{}}
\hline
\rule[-0.2cm]{0pt}{0.7cm} & & & & & & & \\
\hline
\end{tabular}
\par\vspace{0.3cm}
}

% місце для розв'язання: сітка «в клітинку» на всю ширину рядка
\newcommand{\nmtWorkspace}[1]{\par\nopagebreak\vspace{0.2cm}\noindent
\begin{tikzpicture}
\draw[gray!22, step=0.5cm] (0,0) grid ({\linewidth-1pt},#1);
\draw[gray!55] (0,0) rectangle ({\linewidth-1pt},#1);
\node[anchor=north west, inner sep=2pt, xshift=1pt, yshift=-1pt, fill=white, font=\scriptsize, text=gray] at (0,#1) {Місце для розв'язання};
\end{tikzpicture}\par\vspace{0.1cm}}
\begin{document}
\relpenalty=10000 \binoppenalty=10000   % формула в рядку не розривається після «=» чи «+»
% макроси бази
% --- сумісність зі старою базою (діє, якщо тема не має власного визначення) ---
\providecommand{\trigStyle}[1]{#1}
\providecommand{\answerTableSmall}[5]{%
\begingroup\setlength{\tabcolsep}{2pt}\small\nmtSetCell
\begin{tabular}{|*{5}{>{\centering\arraybackslash}m{\nmtcellw}|}}
\hline
\rule[-0.2cm]{0pt}{0.6cm}\textbf{А} & \textbf{Б} & \textbf{В} & \textbf{Г} & \textbf{Д} \\
\hline
\rule[-0.4cm]{0pt}{0.9cm}#1 & \rule[-0.4cm]{0pt}{0.9cm}#2 & \rule[-0.4cm]{0pt}{0.9cm}#3 & \rule[-0.4cm]{0pt}{0.9cm}#4 & \rule[-0.4cm]{0pt}{0.9cm}#5 \\
\hline
\end{tabular}\endgroup}
\providecommand{\matchingLayout}[3]{%
    \noindent
    \begin{minipage}[t]{0.36\textwidth}\vspace{0pt}\raggedright
        #1
    \end{minipage}%
    \hfill
    \begin{minipage}[t]{0.43\textwidth}\vspace{0pt}\raggedright
        #2
    \end{minipage}%
    \hfill
    \begin{minipage}[t]{0.19\textwidth}
        \vspace{0pt}
        \begin{flushright}
        #3
        \end{flushright}
    \end{minipage}}
\providecommand{\matchingLayoutBottom}[2]{%
    \noindent
    \begin{minipage}[t]{0.48\textwidth}\vspace{0pt}\raggedright
        #1
    \end{minipage}%
    \hfill
    \begin{minipage}[t]{0.48\textwidth}\vspace{0pt}\raggedright
        #2
    \end{minipage}
    \vspace{0.5cm}
    \noindent
    \matchingGrid}
\providecommand{\answerListVertical}[5]{%
    \par\nopagebreak\vspace{0.2cm}
    \begin{itemize}[itemsep=0.4cm, leftmargin=1.5cm, labelsep=0.5cm]
        \item[\textbf{А}] #1
        \item[\textbf{Б}] #2
        \item[\textbf{В}] #3
        \item[\textbf{Г}] #4
        \item[\textbf{Д}] #5
    \end{itemize}
    \vspace{0.2cm}}
\setcounter{zad}{0}

\chapterTitle{НМТ 2026 з математики: варіант за 8 червня}

\noindent{\small Оригінальні завдання основної сесії НМТ 2026 з математики (8~червня): 15~завдань з вибором однієї відповіді, 3~на встановлення відповідності й 4~з короткою відповіддю (максимум 32 бали). Під кожним завданням -- місце для розв'язання; відповіді й розв'язки -- окремим файлом.}\par\vspace{0.4cm}

\instructionBox{Завдання 1--15 мають по п'ять варіантів відповіді, з яких лише \textbf{один правильний}. Виберіть правильний, на Вашу думку, варіант відповіді й позначте його в бланку відповідей.}

\begin{samepage}
\noindent\begin{minipage}{\linewidth}
\zadtask{Пляшку об'ємом $750$~мл на $\dfrac{2}{3}$ заповнили соком. Скільки соку залишилося в ній, якщо відлили $150$~мл соку?}
\answerTable{$250$~мл}{$600$~мл}{$300$~мл}{$450$~мл}{$350$~мл}

\nmtWorkspace{4cm}
\end{minipage}
\par\nbvspace{0.35cm}\penalty-20
\par\penalty-20
\end{samepage}

\begin{samepage}
\noindent\begin{minipage}{\linewidth}
\noindent
\begin{minipage}[c]{0.48\textwidth}
\raggedright\zadnum На рисунку зображено графік залежності кількості переглядів відеоролика (у тисячах) від часу (у годинах) з моменту його публікації. Укажіть проміжок часу, протягом якого кількість переглядів зростала найшвидше.
\end{minipage}%
\hfill
\begin{minipage}[c]{0.50\textwidth}
\centering
\begin{nmtfit}\begin{tikzpicture}
\begin{axis}[
    width=8.5cm, height=5.5cm,
    xmin=0, xmax=21,
    ymin=0, ymax=24,
    xtick={0,2,4,6,8,10,12,14,16,18,20},
    ytick={0,4,8,12,16,20,24},
    xlabel={\small Години},
    ylabel={\small Кількість переглядів, тис.},
    grid=both,
    grid style={dashed, gray!40},
    axis lines=left,
    tick label style={font=\scriptsize},
    label style={font=\scriptsize}
]
\addplot[mark=*, thick, red!80!black] coordinates {
    (0,0) (2,1) (4,2) (6,3) (8,6) (10,12) (12,15) (14,16) (16,17) (18,17.5) (20,18)
};
\end{axis}
\end{tikzpicture}\end{nmtfit}
\end{minipage}
\par\nbvspace{0.2cm}
\answerTable{$[0; 4]$}{$[4; 8]$}{$[8; 10]$}{$[12; 16]$}{$[18; 20]$}

\nmtWorkspace{4cm}
\end{minipage}
\par\nbvspace{0.35cm}\penalty-20
\par\penalty-20
\end{samepage}

\begin{samepage}
\noindent\begin{minipage}{\linewidth}
\noindent
\begin{minipage}[c]{0.60\textwidth}
\raggedright\zadnum Ділянка складається з прямокутного трикутника та прямокутника (див. рисунок). За наведеними на рисунку розмірами сторін визначте периметр цієї ділянки.
\end{minipage}%
\hfill
\begin{minipage}[c]{0.38\textwidth}
\centering
\begin{nmtfit}\begin{tikzpicture}[scale=0.3, thick, line join=round]
    \coordinate (A) at (0,0);
    \coordinate (D) at (10,0);
    \coordinate (C) at (10,7);
    \coordinate (B) at (0,7);
    \coordinate (Top) at (6.4, 11.8);
    \draw (A) -- (D) -- (C) -- (Top) -- (B) -- cycle;
    \draw[dashed] (B) -- (C);
    \pic[draw, angle radius=0.3cm] {right angle = A--B--C};
    \pic[draw, angle radius=0.3cm] {right angle = B--A--D};
    \pic[draw, angle radius=0.3cm] {right angle = A--D--C};
    \pic[draw, angle radius=0.3cm] {right angle = D--C--B};
    \pic[draw, angle radius=0.3cm] {right angle = B--Top--C};
    \node[left] at (0,3.5) {$7$ м};
    \node[right] at (10,3.5) {$7$ м};
    \node[above left] at (3.2, 9.4) {$8$ м};
    \node[above right] at (8.2, 9.4) {$6$ м};
\end{tikzpicture}\end{nmtfit}
\end{minipage}
\par\nbvspace{0.2cm}
\answerTable{$28$~м}{$35$~м}{$48$~м}{$38$~м}{$94$~м}

\nmtWorkspace{4cm}
\end{minipage}
\par\nbvspace{0.35cm}\penalty-20
\par\penalty-20
\end{samepage}

\begin{samepage}
\noindent\begin{minipage}{\linewidth}
\zadtask{Позначте число, що є розв'язком системи нерівностей $\begin{cases}x+10>0,\\ 5-x>0.\end{cases}$}
\answerTable{$-5$}{$-10$}{$10$}{$15$}{$5$}

\nmtWorkspace{4cm}
\end{minipage}
\par\nbvspace{0.35cm}\penalty-20
\par\penalty-20
\end{samepage}

\begin{samepage}
\noindent\begin{minipage}{\linewidth}
\noindent
\begin{minipage}[c]{0.62\textwidth}
\raggedright\zadnum На рисунку зображено правильну чотирикутну піраміду $SABCD$, $SO$~--- її висота. Скільки всього є прямих, перпендикулярних до $SO$ та які проходять через дві вершини основи піраміди?\par
\nopagebreak\nbvspace{0.2cm}
\par\noindent\textbf{А}\ \ лише одна
\par\noindent\textbf{Б}\ \ дві
\par\noindent\textbf{В}\ \ чотири
\par\noindent\textbf{Г}\ \ шість
\par\noindent\textbf{Д}\ \ більше шести
\end{minipage}%
\hfill
\begin{minipage}[c]{0.36\textwidth}
\centering
\begin{nmtfit}\begin{tikzpicture}[scale=0.8, thick, line join=round]
    \coordinate (A) at (0,0);
    \coordinate (B) at (1.5,1.2);
    \coordinate (C) at (5,1.2);
    \coordinate (D) at (3.5,0);
    \coordinate (O) at (2.5,0.6);
    \coordinate (S) at (2.5,4.5);
    \draw (A) -- (D) -- (C) -- (S) -- cycle;
    \draw (S) -- (D);
    \draw[dashed] (A) -- (B) -- (C);
    \draw[dashed] (S) -- (B);
    \draw[dashed] (A) -- (C);
    \draw[dashed] (B) -- (D);
    \draw[dashed] (S) -- (O);
    \node[below left] at (A) {$A$};
    \node[above left] at (B) {$B$};
    \node[right] at (C) {$C$};
    \node[below right] at (D) {$D$};
    \node[below] at (O) {$O$};
    \node[above] at (S) {$S$};
\end{tikzpicture}\end{nmtfit}
\end{minipage}
\par\nbvspace{0.3cm}

\nmtWorkspace{4cm}
\end{minipage}
\par\nbvspace{0.35cm}\penalty-20
\par\penalty-20
\end{samepage}

\begin{samepage}
\noindent\begin{minipage}{\linewidth}
\zadtask{$\dfrac{a^2-4b^2}{3a+6b} =$}
\answerTableTall{$\dfrac{a+2b}{2}$}{$\dfrac{a-2b}{3}$}{$3a-6b$}{$\dfrac{3}{a-2b}$}{$3(a+2b)$}

\nmtWorkspace{4cm}
\end{minipage}
\par\nbvspace{0.35cm}\penalty-20
\par\penalty-20
\end{samepage}

\begin{samepage}
\noindent\begin{minipage}{\linewidth}
\zadtask{Якому проміжку належить корінь рівняння $2^{2x-5} = 8$?}
\answerTable{$(-\infty; -2)$}{$[0; 2)$}{$[-2; 0)$}{$[2; 4)$}{$[4; +\infty)$}

\nmtWorkspace{4cm}
\end{minipage}
\par\nbvspace{0.35cm}\penalty-20
\par\penalty-20
\end{samepage}

\begin{samepage}
\noindent\begin{minipage}{\linewidth}
\noindent
\begin{minipage}[c]{0.58\textwidth}
\raggedright\zadnum У прямокутній системі координат на площині зображено точки $K$, $L$, $M$, $N$ і $P$. Через яку із цих точок проходить графік функції $y=3$?
\end{minipage}%
\hfill
\begin{minipage}[c]{0.40\textwidth}
\centering
\begin{nmtfit}\begin{tikzpicture}[scale=0.6, thick]
    \draw[step=1cm, gray!40, thin] (-4,-3) grid (5,5);
    \draw[->, thick] (-4.5,0) -- (5.5,0) node[right] {$x$};
    \draw[->, thick] (0,-3.5) -- (0,5.5) node[above] {$y$};
    \node[below left] at (0,0) {$0$};
    \node[below right] at (1,0) {$1$};
    \node[above left] at (0,1) {$1$};

    \fill (-3,2) circle (3pt) node[above left] {$M$};
    \fill (1,3) circle (3pt) node[above left] {$N$};
    \fill (4,2) circle (3pt) node[above left] {$L$};
    \fill (0,-2) circle (3pt) node[below left] {$K$};
    \fill (3,-3) circle (3pt) node[above right] {$P$};
\end{tikzpicture}\end{nmtfit}
\end{minipage}
\par\nbvspace{0.2cm}
\answerTable{$N$}{$K$}{$L$}{$M$}{$P$}

\nmtWorkspace{4cm}
\end{minipage}
\par\nbvspace{0.35cm}\penalty-20
\par\penalty-20
\end{samepage}

\begin{samepage}
\noindent\begin{minipage}{\linewidth}
\zadtask{Які з наведених тверджень є правильними?\par
\nopagebreak\nbvspace{0.1cm}
\begin{enumerate}[label=\textbf{\Roman*.}, align=left, leftmargin=2.6em, labelsep=0.5em, itemsep=2pt, topsep=2pt]
\item У будь-який ромб можна вписати коло.
\item У будь-яку прямокутну трапецію можна вписати коло.
\item У будь-якому трикутнику бісектриси кутів проходять через центр вписаного в цей трикутник кола.
\end{enumerate}
\nopagebreak\nbvspace{0.15cm}}
\answerTable{лише III}{лише I та II}{лише I та III}{лише II та III}{I, II та III}

\nmtWorkspace{4cm}
\end{minipage}
\par\nbvspace{0.35cm}\penalty-20
\par\penalty-20
\end{samepage}

\begin{samepage}
\noindent\begin{minipage}{\linewidth}
\zadtask{Визначте всі критичні точки функції $f(x) = \dfrac{x^3}{3} - 2x^2 + 3x + 1$.}
\answerTable{$-3; -1$}{$1; 3$}{$3$}{$-3; 1$}{$1$}

\nmtWorkspace{4cm}
\end{minipage}
\par\nbvspace{0.35cm}\penalty-20
\par\penalty-20
\end{samepage}

\begin{samepage}
\noindent\begin{minipage}{\linewidth}
\zadtask{Комп'ютерна програма видаляє у п'ятицифровому числі одну цифру навмання. Яка ймовірність того, що в числі $37281$ буде видалено цифру $1$ або цифру $2$?}
\answerTableTall{$\dfrac{1}{2}$}{$\dfrac{2}{5}$}{$\dfrac{1}{5}$}{$\dfrac{2}{3}$}{$\dfrac{1}{3}$}
\nmtWorkspace{4cm}
\end{minipage}
\par\nbvspace{0.35cm}\penalty-20
\par\penalty-20
\end{samepage}

\begin{samepage}
\noindent\begin{minipage}{\linewidth}
\noindent
\begin{minipage}[c]{0.60\textwidth}
\raggedright\zadnum У прямокутній системі координат у просторі зображено сферу радіуса $4$ (див. рисунок). Точки $M(1; -2; 0)$ і $N$ належать діаметру сфери. Визначте координати точки $N$.
\end{minipage}%
\hfill
\begin{minipage}[c]{0.38\textwidth}
\centering
\begin{nmtfit}\begin{tikzpicture}[scale=0.7, thick, line join=round]
    %
    \draw[fill=yellow!30!orange!20, draw=black] (-2,-1) -- (4,-1) -- (5.5,1.5) -- (-0.5,1.5) -- cycle;
    %
    \draw[->, >=stealth] (2.5,0.5) -- (1,-1.5) node[left] {$x$};
    \draw[->, >=stealth] (2.5,0.5) -- (5,0.5) node[below] {$y$};
    \draw[->, >=stealth] (2.5,0.5) -- (2.5,4.5) node[right] {$z$};
    %
    \coordinate (C) at (1, 2.5); %
    \draw[fill=cyan!15, draw=blue!60!black] (C) circle (1.5);
    \draw[dashed, blue!60!black] (1, 2.5) ellipse (1.5 and 0.4); %
    %
    \coordinate (M) at (1, 1);
    \coordinate (N) at (1, 4);
    \draw[dashed, red!80!black] (M) -- (N);
    \fill (M) circle (2pt) node[left] {$M$};
    \fill (N) circle (2pt) node[left] {$N$};
\end{tikzpicture}\end{nmtfit}
\end{minipage}
\par\nbvspace{0.2cm}
\answerTable{$(1; -2; 4)$}{$(1; -2; 8)$}{$(4; -2; 4)$}{$(5; -2; 0)$}{$(-1; 2; 8)$}

\nmtWorkspace{4cm}
\end{minipage}
\par\nbvspace{0.35cm}\penalty-20
\par\penalty-20
\end{samepage}

\begin{samepage}
\noindent\begin{minipage}{\linewidth}
\zadtask{$\dfrac{(a^3)^4}{a^4 \cdot a^8} =$}
\answerTableTall{$1$}{$a^{\frac{1}{2}}$}{$a^6$}{$0$}{$-1$}

\nmtWorkspace{4cm}
\end{minipage}
\par\nbvspace{0.35cm}\penalty-20
\par\penalty-20
\end{samepage}

\begin{samepage}
\noindent\begin{minipage}{\linewidth}
\noindent
\begin{minipage}[c]{0.58\textwidth}
\raggedright\zadnum На стороні $BC$ прямокутника $ABCD$ вибрано точку $K$ так, що прямі $AK$ і $CD$ перетинаються в точці $M$ (див. рисунок). $MD = 12$~см, $MC : CD = 1 : 3$. Визначте площу прямокутника.
\end{minipage}%
\hfill
\begin{minipage}[c]{0.40\textwidth}
\centering
\begin{nmtfit}\begin{tikzpicture}[scale=0.3, thick, line join=round]
    \coordinate (A) at (0,0);
    \coordinate (D) at (15,0);
    \coordinate (C) at (15,5);
    \coordinate (B) at (0,5);
    \coordinate (K) at (10,5);
    \coordinate (M) at (15,7.5);
    \draw (A) -- (B) -- (C) -- (D) -- cycle;
    \draw (A) -- (M);
    \draw (D) -- (M);
    \fill (K) circle (3pt);
    \fill (M) circle (3pt);
    \node[left] at (A) {$A$};
    \node[above left] at (B) {$B$};
    \node[right] at (C) {$C$};
    \node[right] at (D) {$D$};
    \node[above left] at (K) {$K$};
    \node[above right] at (M) {$M$};
    \pic[draw, angle radius=0.8cm, pic text={$30^\circ$}, angle eccentricity=1.5, font=\scriptsize] {angle = D--A--M};
\end{tikzpicture}\end{nmtfit}
\end{minipage}
\par\nbvspace{0.2cm}
\answerTable{$54\sqrt{3}$}{$81\sqrt{3}$}{$108\sqrt{3}$}{$216$}{$162$}

\nmtWorkspace{4cm}
\end{minipage}
\par\nbvspace{0.35cm}\penalty-20
\par\penalty-20
\end{samepage}

\begin{samepage}
\noindent\begin{minipage}{\linewidth}
\zadtask{Розв'яжіть рівняння $\log_5 x + \log_5(4x) = \log_5 25$.}
\answerTableTall{$-\dfrac{5}{2}; \dfrac{5}{2}$}{$\dfrac{5}{2}$}{$\dfrac{2}{5}$}{$-5$}{$5$}

\nmtWorkspace{4cm}
\end{minipage}
\par\nbvspace{0.35cm}\penalty-20
\par\penalty-20
\end{samepage}

\instructionBox{У завданнях 16--18 до кожного з трьох рядків інформації, позначених цифрами, доберіть один правильний, на Вашу думку, варіант, позначений буквою.}

\begin{samepage}
\noindent\begin{minipage}{\linewidth}
\zadtask{Узгодьте вираз \mbox{(1--3)} із проміжком \mbox{(А--Д)}, якому належить значення цього виразу.\par
\nopagebreak\nbvspace{0.3cm}
\noindent
\matchingLayout{
\matchHead{Вираз}
\matchItem{1}{$\log_{0{,}5} 2$}
\matchItem{2}{$\sqrt{3^2 + (-2)^2}$}
\matchItem{3}{$\left| 2 \cdot \dfrac{2}{7} - 2\dfrac{2}{7} \right|$}
}{
\matchHead{Проміжок}
\matchItem{А}{$[-2; -1)$}
\matchItem{Б}{$[-1; 1)$}
\matchItem{В}{$[1; 2)$}
\matchItem{Г}{$[2; 3)$}
\matchItem{Д}{$[3; 4)$}
}{
\matchingGrid
}
}

\nmtWorkspace{4.5cm}
\end{minipage}
\par\nbvspace{0.35cm}\penalty-20
\par\penalty-20
\end{samepage}

\begin{samepage}
\noindent\begin{minipage}{\linewidth}
\zadtask{Доберіть до функції \mbox{(1--3)} її властивість \mbox{(А--Д)}.\par
\nopagebreak\nbvspace{0.3cm}
\noindent
\matchingLayout{
\matchHead{Функція}
\matchItem{1}{$y = \lg x$}
\matchItem{2}{$y = 0{,}2^x - 1$}
\matchItem{3}{$y = x^2 - 4x + 5$}
}{
\matchHead{Властивість}
\matchItem{А}{є зростаючою}
\matchItem{Б}{є спадною}
\matchItem{В}{не має нулів}
\matchItem{Г}{має лише два нулі}
\matchItem{Д}{є парною}
}{
\matchingGrid
}
}

\nmtWorkspace{4.5cm}
\end{minipage}
\par\nbvspace{0.35cm}\penalty-20
\par\penalty-20
\end{samepage}

\begin{samepage}
\noindent\begin{minipage}{\linewidth}
\zadtask{На рисунку зображено рівні паралелограми $ABCD$ і $AKMD$. $BD \perp AB$. $\angle BAD = 65^\circ$. Установіть відповідність між кутом \mbox{(1--3)} та його градусною мірою \mbox{(А--Д)}.\par
\nopagebreak\nbvspace{0.3cm}
\begin{center}
\begin{nmtfit}\begin{tikzpicture}[scale=0.85, thick, line join=round]
    \coordinate (A) at (0,0);
    \coordinate (D) at (5,0);
    \coordinate (B) at (0.893, 1.915); %
    \coordinate (C) at (5.893, 1.915);
    \coordinate (K) at (0.893, -1.915);
    \coordinate (M) at (5.893, -1.915);

    \draw (-1,0) -- (7,0); %
    \draw (A) -- (B) -- (C) -- (D);
    \draw (A) -- (K) -- (M) -- (D);
    \draw (B) -- (D);

    \pic[draw, angle radius=0.5cm, pic text=$65^\circ$, angle eccentricity=1.7, font=\scriptsize] {angle = D--A--B};
    \pic[draw, angle radius=0.3cm] {right angle = A--B--D};

    \node[above left] at (A) {$A$};
    \node[above left] at (B) {$B$};
    \node[above right] at (C) {$C$};
    \node[above right] at (D) {$D$};
    \node[below left] at (K) {$K$};
    \node[below right] at (M) {$M$};
\end{tikzpicture}\end{nmtfit}
\end{center}
\nopagebreak\nbvspace{0.2cm}
\noindent
\matchingLayout{
\matchHead{Кут}
\matchItem{1}{$\angle ABC$}
\matchItem{2}{$\angle CDM$}
\matchItem{3}{$\angle BDM$}
}{
\matchHead{Градусна міра}
\matchItem{А}{$115^\circ$}
\matchItem{Б}{$125^\circ$}
\matchItem{В}{$130^\circ$}
\matchItem{Г}{$140^\circ$}
\matchItem{Д}{$150^\circ$}
}{
\matchingGrid
}
}

\nmtWorkspace{4.5cm}
\end{minipage}
\par\nbvspace{0.35cm}\penalty-20
\par\penalty-20
\end{samepage}

\instructionBox{Розв'яжіть завдання 19--22. Одержані числові відповіді запишіть у бланку відповідей. Відповідь записуйте лише десятковим дробом, урахувавши положення коми. Знак «мінус» записуйте перед першою цифрою числа.}

\begin{samepage}
\noindent\begin{minipage}{\linewidth}
\zadtask{Знайдіть площу фігури, обмеженої графіком функції $y=3-\sqrt{x}$ та осями координат.}
\nmtAnswerBox

\nmtWorkspace{7cm}
\end{minipage}
\par\nbvspace{0.35cm}\penalty-20
\par\penalty-20
\end{samepage}

\begin{samepage}
\noindent\begin{minipage}{\linewidth}
\zadtask{Михайло планував купити мобільний телефон, чохол до нього та карту пам'яті. Вартість телефона становить $8000$~грн, чохла~--- $600$~грн, а карти пам'яті~--- $800$~грн. У магазині проходить акція: за умови одночасної купівлі телефона та чохла покупець отримує знижку розміром $15\,\%$ від вартості телефона, знижку на чохол у розмірі п'ятої частини від його повної вартості, a карту пам'яті у подарунок. Яку суму грошей (у грн) заплатить Михайло за вибрані ним телефон, чохол та карту пам'яті, якщо скористається цією акцією?}
\nmtAnswerBox

\nmtWorkspace{7cm}
\end{minipage}
\par\nbvspace{0.35cm}\penalty-20
\par\penalty-20
\end{samepage}

\begin{samepage}
\noindent\begin{minipage}{\linewidth}
\noindent
\begin{minipage}[c]{0.60\textwidth}
\raggedright\zadnum На рисунку зображено правильну трикутну призму та циліндр. Об'єм циліндра дорівнює $80\pi$~см$^3$, а площа основи призми дорівнює $12\sqrt{3}$~см$^2$. Радіус основи циліндра дорівнює радіусу кола, вписаного в основу призми. Визначте об'єм (у см$^3$) призми, якщо висота призми відноситься до висоти циліндра як $\sqrt{3} : 1$.
\end{minipage}%
\hfill
\begin{minipage}[c]{0.38\textwidth}
\centering
\begin{nmtfit}\begin{tikzpicture}[scale=0.6, thick, line join=round]
    %
    \coordinate (A1) at (0,0);
    \coordinate (B1) at (2,-1);
    \coordinate (C1) at (3,0.5);
    \coordinate (A2) at (0,4);
    \coordinate (B2) at (2,3);
    \coordinate (C2) at (3,4.5);

    \draw (A1) -- (B1) -- (B2) -- (A2) -- cycle;
    \draw (B1) -- (C1) -- (C2) -- (B2);
    \draw (A2) -- (C2);
    \draw[dashed] (A1) -- (C1);

    %
    \begin{scope}[shift={(6,0)}]
        \draw[dashed] (-1,0) arc (180:0:1 and 0.4);
        \draw (-1,0) arc (180:360:1 and 0.4);
        \draw (-1,0) -- (-1,4);
        \draw (1,0) -- (1,4);
        \draw (0,4) ellipse (1 and 0.4);
    \end{scope}
\end{tikzpicture}\end{nmtfit}
\end{minipage}
\par\nbvspace{0.2cm}
\nmtAnswerBox

\nmtWorkspace{7cm}
\end{minipage}
\par\nbvspace{0.35cm}\penalty-20
\par\penalty-20
\end{samepage}

\begin{samepage}
\noindent\begin{minipage}{\linewidth}
\zadtask{Визначте суму всіх цілих значень $a$, за кожного з яких рівняння $$\frac{9 \cos x - 2a - 5}{\sqrt{\pi^2 - 4x^2}} = 0$$ має хоча б один корінь.}
\nmtAnswerBox

\nmtWorkspace{7cm}
\end{minipage}
\par\nbvspace{0.35cm}\penalty-20
\par\penalty-20
\end{samepage}

\end{document}
